\centerline{\bf \Huge ЕГЭ2024 + ЕГЭотАА}
\centerline{\bf \Huge }

\begin{flushright}
        \scriptsize{\textbf{20:31. Прибыл Годжо Сатору}}
\end{flushright}

\problem{1} Четырёхугольник $ABCD$ вписан в окружность. Угол $ABC$ равен 107°, угол $CAD$ равен 43°. Найдите угол $ABD$. Ответ дайте в градусах.

\problem{2} Даны векторы $\vec{a}(3;0)$ и $\vec{b}(1;3)$. Найдите длину вектора $\vec{a} + 5\vec{b}$.

\problem{3} Найдите объём многогранника, вершинами которого являются вершины $A, B, C, D, B_1$ прямоугольного параллелепипеда $ABCDA_1B_1C_1D_1$, у которого $AB=10$, $BC=4$, $BB_1=9$.

\problem{4*} В аудитории 15 человек и каждый пятый сможет пройти данный отбор. Какова вероятность того, что ВЫ сможете пройти этот отбор?

\problem{5*} Гениальный выпускник сдаёт три экзамена для поступления в ВУЗ мечты. Он мощно готовился, поэтому вероятность сдачи каждого экзамена на 90 баллов или выше равна 0,99. Найдите вероятность того, что выпускник сдаст первый экзамен (по профмату) на 100 баллов, а все остальные экзамены на 1 балл.

\problem{6} Найдите корень уравнения $\sqrt{53-2x} = 7$.

\problem{7} Найдите значение выражения $3\sin\dfrac{13\pi}{12}\cdot \cos\dfrac{13\pi}{12}$.

\problem{8*} Материальная точка движется прямолинейно по закону $x(t) = t^{100} - 99t^{98} + 2024t + 2025$ (где $x$ — расстояние от точки отсчета в метрах, $t$ — время в секундах, измеренное с начала движения). Найдите ее скорость в (м/с) в момент времени $t = 0$ с.

\problem{9} Мотоциклист, движущийся по городу со скоростью $v_0=95$ км/ч, выезжает из него и сразу после выезда начинает разгоняться с постоянным ускорением $a=40$ км/ч$^2$. Расстояние (в км) от мотоциклиста до города вычисляется по формуле $S=v_0t + \dfrac{at^2}{2}$, где $t$ - время в часах, прошедшее после выезда из города. Определите время, прошедшее после выезда мотоциклиста из города, если известно, что за это время он удалился от города на 25 км. Ответ дайте в минутах.

\problem{10*} Одиннадцатиклассник и десятиклассник, работая вместе, решают этот вариант за 40 минут, а один одиннадцатиклассник - за 60 минут. За сколько минут решает этот вариант один десятиклассник?


\problem{11*} Точки $A(-2;1)$ и $B(-1;2)$ принадлежат графику функции $f(x)=x^3-zx+v$. Найдите f(22).

\problem{12} Найдите точку максимума функции $y=8\cdot \ln{(x-5)} - 8x + 9$.

\problem{13} а) Решите уравнение 

$\cos{2x}-\sin{(x-\pi)} - 1 = 0$.

б) Укажите корни этого уравнения, принадлежащие отрезку $[-\frac{7\pi}{2};-2\pi]$.

\problem{14*} В пирамиде $SBCD$ каждое ребро равно 3. На ребре $SB$ взята точка $A$ так, что $SA:AB=1:2$. Найдите радиус сферы, описанной около пирамиды $SACD$.


\problem{15} Решите неравенство

$ \dfrac{2^x+8}{2^x-8} + \dfrac{2^x-8}{2^x+8} \geq \dfrac{2^{(x+4)}+96}{4^x-64}$

\problem{16*} В воскресенье, 20 октября, на Пагоде семь гопников остановили Анджура Аркадьевича. После недлительных переговоров АА предложил гопникам интересную схему заёма денежных средств. Каждый понедельник долг увеличивается на 20\%, и со вторника по субботу необходимо выплатить одним платежом часть долга. Гопники три раза переводили одну и ту же сумму на оффшорный счет АА и полностью погасили свой долг. Сколько рублей они изначально занимали у АА, если им пришлось вернуть на 77200 рублей больше суммы, взятой в долг?

\problem{17} Пятиугольник $ABCDE$ вписан в окружность. Известно, что $AB=CD=3$, $BC=DE=4$.

а) Докажите, что $AC=CE$

б) Найдите длину диагонали $BE$, если $AD=6$.

\problem{18*} При каких значениях параметра $z$ уравнение

$z^4x + z^2 +(2+\sqrt{2})z + 2\sqrt{2}=z^2(z+\sqrt{2})+4x$

имеет бесконечно много корней?

\problem{19*} Три числа, сумма которых равна 28, образуют геометрическую прогрессию. Если к первому числу прибавить 3, ко второму числу прибавить 1, а от третьего числа отнять 5, то полученные числа образуют арифметическую прогрессию. Найдите эти числа.


